\documentclass[aspectratio=169]{beamer}
\usepackage[spanish]{babel}
\usepackage[utf8]{inputenc}
\usepackage{tikz}
\usepackage{booktabs}
\usepackage{amsmath}
\usepackage{amssymb}
\usepackage{svg}
\usepackage{graphicx}
\graphicspath{{imagenes/}}
\usetheme{Berkeley}
\usecolortheme{default}

\title[LCL Conmutado]{Diseño, modelado e implementación de un circuito conmutado para la protección por sobre-corriente}
\subtitle{Trabajo de Fin de Grado}
\author[Á. Villalta]{Álvaro Villalta Vázquez}

\institute[UC3M]{
  Grado en Ingeniería Electrónica Industrial y Automática\\
  Universidad Carlos III de Madrid\\[4pt]
  \textit{Tutora:} Cristina Fernández Herrero
}

\date[2025-2026]{Leganés, 2026}
\logo{\includegraphics[height=1.2cm]{UC3M (1).png}}

\AtBeginSection[]
{
  \begin{frame}
    \frametitle{Índice}
    \tableofcontents[currentsection]
  \end{frame}
}

\begin{document}

%--------------------------------------------------------------
% 1. PORTADA
%--------------------------------------------------------------
\frame{\titlepage}

%--------------------------------------------------------------
% 2. ÍNDICE GENERAL
%--------------------------------------------------------------
\begin{frame}
\frametitle{Índice}
\tableofcontents
\end{frame}

%==============================================================
\section{Introducción y Motivación}
%==============================================================

%--------------------------------------------------------------
% 3. CONTEXTO
%--------------------------------------------------------------
\begin{frame}
\frametitle{Contexto: protección en buses aeroespaciales}

Los sistemas de distribución de potencia en satélites requieren protección fiable ante sobrecorrientes.

\pause
\begin{block}{¿Qué es un LCL?}
Un \textbf{Latching Current Limiter (LCL)} es un dispositivo que:
\begin{itemize}
  \item Limita la corriente ante una sobrecarga
  \pause
  \item Mantiene la carga conectada durante un tiempo prefijado (\textit{trip-off time})
  \pause
  \item Se \textbf{enclava} en estado abierto si la falta persiste
\end{itemize}
\end{block}

\pause
\begin{alertblock}{Ventaja clave frente a fusibles}
Permite rearmado posterior (por telecomando), sin sustitución de componentes.
\end{alertblock}

\end{frame}

%--------------------------------------------------------------
% 4. PROBLEMÁTICA P-MOS
%--------------------------------------------------------------
\begin{frame}
\frametitle{Problemática: topología clásica P-MOS}

Las topologías tradicionales emplean MOSFET de canal P en \textbf{región lineal}:

\pause
\begin{itemize}
  \item \textbf{Ventaja:} control de puerta sencillo (tensión $V_G < V_{bus}$)
  \pause
  \item \textbf{Problema 1:} $R_{DS,ON}$ del P-MOS es $\approx 3\times$ mayor que el N-MOS equivalente
  \pause
  \item \textbf{Problema 2:} Durante la limitación, el transistor disipa:
  \[
    P_{lineal} = V_{DS} \cdot I_{lim} \approx V_{bus} \cdot I_{lim}
  \]
  \pause
  \item \textbf{Consecuencia:} A 100 V, el estrés térmico es \alert{inviable} con encapsulados estándar
\end{itemize}

\pause
\begin{examples}
A $V_{bus} = 100$ V e $I_{lim} = 4$ A la potencia instantánea es \textbf{400 W} en el semiconductor.
\end{examples}

\end{frame}

%--------------------------------------------------------------
% 5. SOLUCIÓN PROPUESTA
%--------------------------------------------------------------
\begin{frame}
\frametitle{Solución propuesta: N-MOS en conmutación}

\begin{block}{Propuesta}
Limitador en \textbf{topología reductora (Buck)} con N-MOS en high-side y control por histéresis.
\end{block}

\pause
\begin{itemize}
  \item El transistor \textbf{conmuta} a alta frecuencia durante la limitación
  \pause
  \item Las pérdidas se reducen drásticamente: conducción + conmutación en vez de disipación lineal
  \pause
  \item Permite \textit{trip-off times} de hasta \textbf{1,5 ms} a 100 V con semiconductores estándar
  \pause
  \item Topología basada en el trabajo de \textit{López et al. (SAAEI 2024)}
\end{itemize}

\pause
\begin{alertblock}{Reto del N-MOS en high-side}
La excitación de puerta requiere $V_G > V_{bus}$: se emplea un \textbf{convertidor CC-CC aislado flotante}.
\end{alertblock}

\end{frame}

%--------------------------------------------------------------
% 6. OBJETIVOS
%--------------------------------------------------------------
\begin{frame}
\frametitle{Objetivos del TFG}

\textbf{Objetivo principal:}
\begin{block}{}
Diseñar, simular y validar experimentalmente un LCL en conmutación N-MOS, demostrando la mejora frente a topologías lineales con P-MOS.
\end{block}

\pause
\textbf{Objetivos específicos:}
\begin{itemize}
  \item[$\checkmark$] Analizar el estado del arte de los LCL y justificar la transición P-MOS $\to$ N-MOS
  \pause
  \item[$\checkmark$] Dimensionar teóricamente la topología: inductancia, histéresis, temporizador
  \pause
  \item[$\checkmark$] Modelar el circuito en \textbf{PSIM} con parásitos de componentes reales
  \pause
  \item[$\checkmark$] Construir y ensayar una \textbf{PCB} en laboratorio
  \pause
  \item[$\checkmark$] Definir metodología de \textbf{escalado} a otros buses y corrientes
\end{itemize}

\end{frame}

%==============================================================
\section{Estado del Arte}
%==============================================================

%--------------------------------------------------------------
% 7. EVOLUCIÓN DE LA PROTECCIÓN
%--------------------------------------------------------------
\begin{frame}
\frametitle{Evolución de los sistemas de protección}

\begin{enumerate}
  \item \textbf{Fusibles} (1864): respuesta inversa corriente-tiempo, no rearmables
  \pause
  \item \textbf{Disyuntores} (1924): rearmables pero con arco DC y desgaste mecánico
  \pause
  \item \textbf{SSCB / SSPC} (estado sólido): sin partes móviles, respuesta en $\mu$s
  \pause
  \item \textbf{LCL}: combina limitación temporal + enclavamiento + rearmado remoto
\end{enumerate}

\pause
\begin{block}{Marco normativo}
La ECSS define el LCL como dispositivo de protección conmutable situado entre fuente y carga que desconecta tras un límite de sobrecorriente durante un tiempo determinado.
\end{block}

\end{frame}

%--------------------------------------------------------------
% 8. P-MOS vs N-MOS
%--------------------------------------------------------------
\begin{frame}
\frametitle{P-MOS vs N-MOS: fundamento físico}

La superioridad del canal N radica en la \textbf{movilidad de portadores}:

\pause
\[
  \mu_e \approx 1350\ \text{cm}^2/\text{V·s} \quad \gg \quad \mu_p \approx 480\ \text{cm}^2/\text{V·s}
\]

\pause
\begin{block}{Consecuencia en $R_{DS,ON}$}
\[
  R_{DS,ON} \propto \frac{1}{\sigma} = \frac{1}{q(n\mu_n + p\mu_p)}
\]
Un N-MOS tiene $R_{DS,ON} \approx 3\times$ menor que un P-MOS equivalente en área de silicio.
\end{block}

\pause
\begin{itemize}
  \item Menor $R_{DS,ON}$ $\Rightarrow$ menos pérdidas en conducción
  \pause
  \item Menor área de silicio $\Rightarrow$ menores capacidades parásitas $\Rightarrow$ mayor velocidad
\end{itemize}

\end{frame}

%--------------------------------------------------------------
% 9. LINEAL vs CONMUTADO
%--------------------------------------------------------------
\begin{frame}
\frametitle{Limitación lineal vs. conmutada}

\begin{columns}
\column{0.48\textwidth}
\begin{alertblock}{Modo lineal (P-MOS)}
\begin{itemize}
  \item MOSFET en zona activa
  \item $P_{dis} = V_{DS} \cdot I_{lim}$
  \item Estrés térmico extremo
  \item Inviable a 100 V
\end{itemize}
\end{alertblock}

\pause
\column{0.48\textwidth}
\begin{block}{Modo conmutado (N-MOS)}
\begin{itemize}
  \item MOSFET alterna ON/OFF
  \item Pérdidas = conducción + switching
  \item Eficiencia $\approx 96{,}5\%$
  \item Escalable a 100 V y más
\end{itemize}
\end{block}
\end{columns}

\pause
\vspace{6pt}
\begin{examples}
El modo conmutado permite \textit{trip-off times} del orden de milisegundos sin riesgo de destrucción del semiconductor, incluso a tensiones de bus elevadas.
\end{examples}

\end{frame}

%==============================================================
\section{Diseño Teórico}
%==============================================================

%--------------------------------------------------------------
% 10. ESPECIFICACIONES
%--------------------------------------------------------------
\begin{frame}
\frametitle{Especificaciones de diseño}

\begin{block}{Bus objetivo: 100 V (aeroespacial)}
Sin clase ECSS estandarizada para este voltaje $\Rightarrow$ diseño personalizado.
\end{block}

\pause
\begin{table}
\centering
\begin{tabular}{ll}
\toprule
\textbf{Parámetro} & \textbf{Valor} \\
\midrule
Tensión de bus ($V_{bus}$) & 100 V \\
\pause
Corriente de limitación ($I_{lim}$) & 3,7 A \\
Banda de histéresis ($\Delta i$) & 0,8 A (3,7–4,5 A) \\
\pause
Tiempo de trip-off ($t_{trip}$) & 1,5 ms \\
\pause
Frecuencia de conmutación ($f_{sw}$) & $\approx$ 615 kHz \\
\bottomrule
\end{tabular}
\end{table}

\pause
\begin{alertblock}{Criterio de diseño}
Se toma el \textit{trip-off} más restrictivo de la clase 1 de 28 V, dada la mayor potencia disipada a 100 V.
\end{alertblock}

\end{frame}

%--------------------------------------------------------------
% 11. ARQUITECTURA
%--------------------------------------------------------------
\begin{frame}
\frametitle{Arquitectura del LCL conmutado}

El sistema se estructura en \textbf{cuatro bloques} principales:

\pause
\begin{enumerate}
  \item \textbf{Etapa de potencia:} Buck N-MOS high-side + diodo Schottky SiC + bobina L
  \pause
  \item \textbf{Sensado y acondicionamiento:} $R_{shunt}$ + AO en cascada (ganancia $G = 25$ V/V)
  \pause
  \item \textbf{Control por histéresis:} comparadores LM319 + biestable RS + AND gate
  \pause
  \item \textbf{Temporizador asimétrico:} red RC con diodo BAT54S + LM393 para memoria térmica
\end{enumerate}

\pause
\begin{block}{Alimentación flotante}
Convertidor CC-CC aislado TMA 1515S que provee +15 V respecto a la fuente del MOSFET, necesario para excitar el N-MOS en high-side con ciclo de trabajo unitario.
\end{block}

\end{frame}

%--------------------------------------------------------------
% 12. DIMENSIONAMIENTO CLAVE
%--------------------------------------------------------------
\begin{frame}
\frametitle{Dimensionamiento: parámetros críticos}

\textbf{Resistencia de sensado} ($R_{shunt} = 22\ \text{m}\Omega$):

\pause
\[
  V_{sense} = I_{lim} \cdot R_{shunt} \cdot G = 3{,}7 \times 0{,}022 \times 25 \approx 2{,}04\ \text{V}
\]

\pause
\vspace{4pt}
\textbf{Inductancia crítica} (criterio ECSS):
\[
  L = \frac{V_{bus} \cdot D(1-D)}{f_{sw} \cdot \Delta i} \quad \Rightarrow \quad L = 68\ \mu\text{H}
\]

\pause
\vspace{4pt}
\textbf{Temporizador} ($R_7 = 26{,}1\ \text{k}\Omega$, $C_{t1} = 100\ \text{nF}$):
\[
  t_{trip} = -R_7 C_{t1} \cdot \ln\!\left(1 - \frac{V_{TimON}}{5\,\text{V} - V_F}\right) = 1{,}5\ \text{ms}
\]

\pause
\begin{block}{Memoria térmica}
La red RC asimétrica acumula carga en fallos sucesivos, reduciendo el trip-off time en cada fallo posterior.
\end{block}

\end{frame}

%==============================================================
\section{Selección de Componentes y Simulación}
%==============================================================

%--------------------------------------------------------------
% 13. COMPONENTES PRINCIPALES
%--------------------------------------------------------------
\begin{frame}
\frametitle{Componentes principales seleccionados}

\begin{itemize}
  \item \textbf{MOSFET:} IRFB4020PBF — N-CH, 200 V, $R_{DS,ON} = 100\ \text{m}\Omega$, $Q_g = 18\ \text{nC}$
  \pause
  \item \textbf{Gate Driver:} UCC27511AQDBVRQ1 — 4 A source / 8 A sink, $t_{pd} = 23\ \text{ns}$
  \pause
  \item \textbf{Diodo libre conducción:} STPSC6H065B (SiC) — 650 V, 6 A, $Q_{rr}$ mínima
  \pause
  \item \textbf{Inductor:} WE-HCI 74435586800 — 68 µH, $I_{sat} = 7{,}5$ A, $R_{DC} = 27{,}3\ \text{m}\Omega$
  \pause
  \item \textbf{Comparadores histéresis:} LM319 — alta velocidad, salida colector abierto, 80 ns
  \pause
  \item \textbf{Comparadores temporizador:} LM393 — propósito general, baja corriente
  \pause
  \item \textbf{Alimentación flotante:} TMA 1515S + L78L05 — +15 V y +5 V aislados
\end{itemize}

\end{frame}

%--------------------------------------------------------------
% 14. SIMULACIÓN PSIM
%--------------------------------------------------------------
\begin{frame}
\frametitle{Simulación en PSIM: resultados}

El modelo incorpora \textbf{parásitos reales}: tiempos de propagación, ESR, capacidades, etc.

\pause
\begin{block}{Resultado ante sobrecorriente permanente}
\begin{itemize}
  \item Corriente limitada entre \textbf{3,7 A y 4,5 A}
  \item Enclavamiento a los \textbf{1,496 ms}
  \item Frecuencia de conmutación: $\approx 615\ \text{kHz}$
\end{itemize}
\end{block}

\pause
\begin{block}{Validación ante fallos sucesivos}
La memoria térmica del temporizador reduce el tiempo de enclavamiento en cada fallo consecutivo, protegiendo el semiconductor frente a disparos repetidos.
\end{block}

\pause
\begin{examples}
La simulación confirma el diseño teórico con tolerancia muy reducida en límites de histéresis y temporización.
\end{examples}

\end{frame}

%==============================================================
\section{PCB y Validación Experimental}
%==============================================================

%--------------------------------------------------------------
% 15. DISEÑO PCB
%--------------------------------------------------------------
\begin{frame}
\frametitle{Diseño de la PCB en KiCad}

La PCB se diseñó siguiendo criterios industriales y espaciales:

\pause
\begin{itemize}
  \item \textbf{Segregación de masas:} plano flotante de potencia separado del plano de control
  \pause
  \item \textbf{Pistas de potencia:} dimensionadas según norma IPC-2221 para $\leq 10^\circ$C de elevación
  \pause
  \item \textbf{Condensadores de desacoplo:} posicionados estratégicamente junto a cada integrado
  \pause
  \item \textbf{Trazas mínimas:} agrupación del lazo de control para minimizar inductancias parásitas
\end{itemize}

\pause
\begin{alertblock}{Problema detectado}
Error de asociación entre huella SOIC-8 y encapsulado real VSSOP-8 del biestable RS \texttt{SN74LVC2G02DCUR}.
\end{alertblock}

\end{frame}

%--------------------------------------------------------------
% 16. MONTAJE EN ARAÑA Y ENSAYOS
%--------------------------------------------------------------
\begin{frame}
\frametitle{Montaje y ensayos experimentales}

\begin{block}{Montaje en araña (\textit{dead-bug wiring})}
El error de huella obligó a soldar el integrado en araña, aumentando la inductancia parásita serie e inestabilizando el circuito de control.
\end{block}

\pause
\textbf{Ensayos realizados:}
\begin{itemize}
  \item Excitación del N-MOS en high-side mediante \textbf{PWM externa} a $\approx 508\ \text{kHz}$
  \pause
  \item Validación de la etapa de potencia a \textbf{15 V} (limitación de fuente disponible)
  \pause
  \item Medición de eficiencia: $\eta = 89{,}2\%$ a $V_{bus} = 12{,}2$ V vs $96{,}5\%$ teórico a 100 V
  \pause
  \item Verificación del \textbf{filtrado inductivo}: bobina convierte onda cuadrada a señal continua
\end{itemize}

\pause
\begin{alertblock}{Limitación}
El circuito de histéresis y el enclavamiento no pudieron validarse por ruido en el biestable y falta de cargas de potencia adecuadas.
\end{alertblock}

\end{frame}

%--------------------------------------------------------------
% 17. FORMAS DE ONDA
%--------------------------------------------------------------
\begin{frame}
\frametitle{Resultados del osciloscopio}

\textbf{Ensayo a $V_{in} = 14$ V, $D \approx 0{,}65$, $f_{sw} \approx 508$ kHz:}

\pause
\begin{itemize}
  \item $V_{out} = 8{,}37$ V (predicho por $V_{out} = D \cdot V_{in} - \text{caídas parásitas}$)
  \pause
  \item Nodo pre-bobina: onda cuadrada 14 V pico a pico con ligero \textit{ringing}
  \pause
  \item Nodo post-bobina: señal continua $\approx 8{,}37$ V con ruido residual mínimo
  \pause
  \item Eficiencia a $V_{in} = 12{,}2$ V: $\eta = \frac{V_{out} \cdot I_{out}}{V_{in} \cdot I_{in}} = 89{,}2\%$
\end{itemize}

\pause
\begin{block}{Conclusión de la etapa de potencia}
El MOSFET conmuta correctamente en high-side y la bobina filtra adecuadamente a alta frecuencia, validando el núcleo del convertidor Buck.
\end{block}

\end{frame}

%==============================================================
\section{Metodología de Escalado}
%==============================================================

%--------------------------------------------------------------
% 18. ESCALADO
%--------------------------------------------------------------
\begin{frame}
\frametitle{Metodología de escalado del diseño}

El diseño es adaptable a otros buses y corrientes cambiando un mínimo de componentes:

\pause
\begin{block}{1. Adaptar corriente de limitación $I_{lim,nuevo}$}
\[
  V_{TimON} = I_{lim,nuevo} \cdot R_{shunt} \cdot G
  \qquad
  \frac{R_1}{R_2} = \frac{5\,\text{V}}{V_{TimON}} - 1
\]
\end{block}

\pause
\begin{block}{2. Adaptar trip-off time}
\[
  R_7 = \frac{-t_{trip,nuevo}}{C_{t1} \cdot \ln\!\left(1 - \dfrac{V_{TimON}}{5\,\text{V} - V_F}\right)}
\]
\end{block}

\pause
\begin{block}{3. Verificar componentes críticos}
$R_{shunt}$ (potencia disipada), inductor (corriente de saturación $I_{sat}$), MOSFET (estrés térmico).
\end{block}

\end{frame}

%==============================================================
\section{Conclusiones y Trabajo Futuro}
%==============================================================

%--------------------------------------------------------------
% 19. CONCLUSIONES
%--------------------------------------------------------------
\begin{frame}
\frametitle{Conclusiones}

\begin{enumerate}
  \item \textbf{Diseño teórico completado:} parámetros dimensionados siguiendo criterios ECSS para bus de 100 V.
  \pause
  \item \textbf{Simulación exitosa en PSIM:} limitación entre 3,7–4,5 A durante 1,496 ms, con enclavamiento y memoria térmica validados.
  \pause
  \item \textbf{PCB fabricada:} diseño siguiendo IPC-2221 con segregación de masas y correcta capacidad de corriente en pistas.
  \pause
  \item \textbf{Etapa de potencia validada:} conmutación del N-MOS en high-side y filtrado inductivo a $\approx 508$ kHz comprobados.
  \pause
  \item \textbf{La topología N-MOS conmutada} demuestra eficiencia térmica superior a los sistemas lineales P-MOS, justificando el esfuerzo adicional en la excitación del transistor.
\end{enumerate}

\pause
\begin{alertblock}{Pendiente de validación}
El lazo de histéresis y el enclavamiento deben ser verificados en una segunda iteración del prototipo.
\end{alertblock}

\end{frame}

%--------------------------------------------------------------
% 20. TRABAJO FUTURO
%--------------------------------------------------------------
\begin{frame}
\frametitle{Trabajo futuro}

\begin{itemize}
  \item \textbf{Corrección de la PCB:} sustituir huella SOIC-8 por VSSOP-8 del biestable RS y reestructurar la zona de control.
  \pause
  \item \textbf{Reconfiguración de histéresis:} adaptar potenciómetros para permitir corrientes de limitación inferiores (más seguras en laboratorio).
  \pause
  \item \textbf{Validación a 100 V:} ensayar con cargas y fuentes adecuadas para verificar limitación, enclavamiento y fallos sucesivos.
  \pause
  \item \textbf{Componentes EEE (certificación espacial):} reemplazar por equivalentes rad-hard, incluyendo MOSFET de SiC.
  \pause
  \item \textbf{Optimización avanzada:} digitalización del control (FPGA/µC certificado), topologías de próxima generación (relé en paralelo con SSCB).
\end{itemize}

\pause
\begin{examples}
La topología desarrollada es directamente escalable a materiales WBG (SiC, GaN) sin cambios en el circuito de control.
\end{examples}

\end{frame}

%--------------------------------------------------------------
% 21. CIERRE
%--------------------------------------------------------------
\begin{frame}
\frametitle{}

\vspace{1.5cm}
\begin{center}
{\Large \textbf{Gracias por vuestra atención}}

\vspace{1cm}
{\large Álvaro Villalta Vázquez}

\vspace{0.4cm}
{\normalsize Grado en Ingeniería Electrónica Industrial y Automática}\\
{\normalsize Universidad Carlos III de Madrid — 2026}

\vspace{1cm}
\textit{¿Preguntas?}
\end{center}

\end{frame}

\end{document}